%% kapa: exam=1 round=32 year=2021 period=2 form=A subject=appraisal_law subject_ko=감정평가관계법규 seq=17 number=17
%% source: Q-Net 감정평가사 기출문제 (cst00309) artlSeq=5208397 file=extracted/1008/5208397/A형/2021년도 제32회 감정평가사 1차 2교시 A형.pdf
%% conversion: kordoc 4.24.51 -> tools/gen1.py (auto); review=proofread-sonnet; flags=none
\begin{question}{17}
감정평가 및 감정평가사에 관한 법령상 감정평가사에 관한 설명으로 옳지 않은 것은?
\begin{choices}
  \item 감정평가사 결격사유는 감정평가사 등록의 거부사유와 취소사유가 된다.
  \item 등록한 감정평가사는 5년마다 그 등록을 갱신하여야 한다.
  \item 등록한 감정평가사가 징계로 감정평가사 자격이 취소된 후 5년이 지나지 아니한 경우 국토교통부장관은 그 등록을 취소할 수 있다.
  \item 감정평가사는 감정평가업을 하기 위하여 1개의 사무소만을 설치할 수 있다.
  \item 부정한 방법으로 감정평가사 자격을 받은 이유로 그 자격이 취소된 후 3년이 지나지 아니한 사람은 감정평가사가 될 수 없다.
\end{choices}
\end{question}
\answer{3}
